\thispagestyle{empty}
\begin{tikzpicture}[remember picture,overlay]
\begin{scope}[shift={(current page.south west)},x=1mm,y=1mm]
\fill[otNavy](0,0)rectangle(210,297);\fill[otNavyDark](0,0)rectangle(210,88);
\draw[white,opacity=.055,step=7mm](0,0)grid(210,297);
\node[anchor=north west,inner sep=0pt]at(23,278){\includegraphics[width=57mm]{images/logo-blanc.png}};
\node[anchor=west,text=otGreen,font=\otSemi\fontsize{12}{15}\selectfont]at(24,235){MPSI\quad /\quad PHYSIQUE};
\node[anchor=west,text=white,font=\sffamily\bfseries\fontsize{52}{56}\selectfont]at(22,211){Physique};
\node[anchor=north west,text=white,text width=158mm,font=\otLight\fontsize{31}{38}\selectfont]at(24,191){Filtrage linéaire};
\fill[otGreen](24,149)rectangle(54,150.4);
\node[anchor=north west,text=white,text width=153mm,font=\sffamily\fontsize{11}{16}\selectfont]at(24,139){Cours développé et schémas vectoriels\par 24 exercices et corrigés expliqués\par 3 défis, TP, activité numérique, DM et DS\par Parcours sur des supports authentiques};
\draw[white,opacity=.6,-{Stealth}] (24,46)--(185,46);
\draw[white,opacity=.6,-{Stealth}] (28,17)--(28,81);
\draw[otGreen,line width=1.8pt] plot[domain=0:6,samples=180] ({30+24*\x},{46+27*cos(deg(3*\x))*exp(-.18*\x)});
\node[anchor=west,text=white,font=\otSemi\fontsize{9}{12}\selectfont]at(24,8){CHAPITRE 07\quad /\quad COLLECTION OPTITAB\quad /\quad OCTOBRE 2026};
\end{scope}\end{tikzpicture}\null

\Page{Collection OptiTAB}{Physique MPSI}{titre}
\thispagestyle{mpsiFront}
\vspace{8mm}\begin{center}\includegraphics[width=63mm]{images/logo-couleur.png}\par\vspace{15mm}
{\otSemi\color{otGreenDark}CLASSE PRÉPARATOIRE SCIENTIFIQUE}\par\vspace{5mm}
{\otLight\fontsize{23}{29}\selectfont Cours et exercices corrigés}\par\vspace{12mm}
{\fontsize{18}{24}\selectfont Filtrage linéaire}\par\vspace{8mm}
Chapitre 07\par 24 exercices progressifs ; 3 défis avancés\par TP, numérique, DM et DS avec barèmes\par\vspace{12mm}
Collection OptiTAB\par\href{https://www.optitab.net}{optitab.net}\end{center}
\vfill\source{Édition de travail du 25 octobre 2026. Énoncés OptiTAB originaux, figures vectorielles et corrigés indépendants. Les supports extérieurs sont identifiés à la fin du chapitre.}


\Page{Pour ouvrir le chapitre}{Avant-propos et mode d’emploi}{avant}
\thispagestyle{mpsiFront}
Ce chapitre prolonge le manuel de physique MPSI dans la charte de la collection. Il traite les spectres de signaux périodiques, les fonctions de transfert, les diagrammes de Bode et le choix d’un filtre selon un cahier des charges. Le programme de référence est l’\href{https://www.education.gouv.fr/bo/21/Special1/ESRS2035779A.htm}{arrêté du 5 janvier 2021}, annexe 2, rubrique 1.5.
\Sub{Comprendre, chercher, vérifier}
Lire un résultat ne suffit pas à se l’approprier. Redessiner le système, orienter les grandeurs, retrouver les équations puis expliquer leurs limites : ces gestes relient le cours aux problèmes. Une réponse numérique doit comporter une unité, un nombre raisonnable de chiffres et une interprétation.
\begin{loi}{Un parcours en quatre étapes}
1. Faire le diagnostic p.\,\pageref{diagnostic} et reprendre les prérequis fragiles.\par
2. Travailler le cours p.\,\pageref{cours-1} et refaire les exemples.\par
3. Chercher les exercices p.\,\pageref{td1}, puis les défis et évaluations.\par
4. Corriger, expliquer l’erreur puis refaire une question sans solution.
\end{loi}
\Sub{Utiliser les corrigés}
Les solutions commencent p.\,\pageref{corr-diag}. Avant de les lire, utiliser les indications p.\,\pageref{indices}. Le barème du DM et du DS sert à distinguer mise en équation, calcul et interprétation ; ce n’est pas un barème officiel de concours.
\Sub{Travailler sur de vrais sujets}
Les pages p.\,\pageref{sources-1} et suivantes donnent les liens, les questions ciblées et des repères de correction de supports réellement publiés. Conserver leurs données et figures d’origine. Les énoncés OptiTAB sont des créations pédagogiques, clairement distinguées de ces documents.
\begin{methode}{Une fiche de progrès}
Après chaque exercice, noter S (réussi seul), A (réussi avec aide) ou R (à reprendre). Choisir une question à refaire deux jours plus tard. Une solution comprise doit pouvoir être reconstruite, pas seulement reconnue.
\end{methode}
\source{© OptiTAB, 2026. Les séries de données créées pour le TP et le numérique sont signalées comme simulations. Les références des auteurs extérieurs restent attachées aux documents d’origine.}


\Page{Collection OptiTAB}{Sommaire}{sommaire}
\thispagestyle{mpsiFront}
\begingroup\setlength{\parskip}{0pt}\fontsize{10}{13.5}\selectfont
\source{MPSI · Physique · Plan de l’ouvrage}
\tocpart{1}{Ondes et signaux}
\planitem{01}{Formation des images et optique géométrique}{1.1}
\planitem{02}{Signaux électriques dans l'ARQS}{1.2}
\planitem{03}{Circuits linéaires du premier ordre}{1.3}
\planitem{04}{L'oscillateur harmonique}{1.4}
\planitem{05}{Oscillations libres amorties}{1.4}
\planitem{06}{Régime sinusoïdal forcé et résonance}{1.4}
\planitem{07}{Filtrage linéaire}{1.5}
\planitem{08}{Propagation d'un signal}{1.6}
\planitem{09}{Superposition et interférences}{1.6}
\planitem{10}{Ondes stationnaires et corde vibrante}{1.6}
\tocpart{2}{Mouvements et interactions}
\planitem{11}{Décrire le mouvement d'un point}{2.1}
\planitem{12}{Lois de Newton et dynamique}{2.2}
\planitem{13}{Énergie du point matériel}{2.3}
\planitem{14}{Particules chargées dans des champs uniformes}{2.4}
\planitem{15}{Moment cinétique}{2.5}
\planitem{16}{Forces centrales et satellites}{2.6}
\planitem{17}{Mouvement d'un solide}{2.7}
\tocpart{3}{L'énergie : conversions et transferts}
\planitem{18}{Décrire un système à l'équilibre}{3.1}
\planitem{19}{Transformations, travail et transfert thermique}{3.2}
\planitem{20}{Premier principe et bilans d'énergie}{3.3}
\planitem{21}{Deuxième principe et bilans d'entropie}{3.4}
\planitem{22}{Machines thermiques}{3.5}
\tocpart{4}{Induction et physique quantique}
\planitem{23}{Champ magnétique}{1.7.1}
\planitem{24}{Actions d'un champ magnétique}{1.7.2}
\planitem{25}{Induction dans un circuit fixe}{1.7.3 et 1.7.4}
\planitem{26}{Induction et circuits mobiles}{1.7.5}
\planitem{27}{Introduction à la physique quantique}{1.8}
\tocpart{5}{Outils transversaux}
\planitem{28}{Mesures, incertitudes et méthodes expérimentales}{Formation expérimentale}

\endgroup
\vspace{2mm}\hrule\vspace{2mm}
\tocresource{Parcours et ressources du chapitre 07}{ressources}
\source{La pagination concerne ce fascicule. Les autres chapitres sont des volumes séparés ou en préparation ; aucune page fictive n’est attribuée.}


\Page{Collection OptiTAB}{Parcours et ressources}{ressources}
\thispagestyle{mpsiFront}
\tocresource{Diagnostic}{diagnostic}
\tocresource{Cours et méthodes}{cours-1}
\tocresource{Travaux pratiques}{tp}
\tocresource{Activité numérique}{numerique}
\tocresource{24 exercices progressifs}{td1}
\tocresource{Trois défis avancés}{defi1}
\tocresource{Devoir maison}{dm}
\tocresource{Devoir surveillé de 2 heures}{ds1}
\tocresource{Indications progressives}{indices}
\tocresource{Corrigés du diagnostic}{corr-diag}
\tocresource{Corrigés des exercices}{corr-td1}
\tocresource{Corrigés des défis}{corr-defi1}
\tocresource{Corrigé et barème du DM}{corr-dm}
\tocresource{Corrigé et barème du DS}{corr-ds1}
\tocresource{Supports authentiques et références}{sources-1}
\begin{docu}{Mon objectif de travail}\checkline{Je refais une démonstration.}\checkline{Je réussis une application sans aide.}\checkline{Je justifie le choix d’un modèle.}\checkline{Je corrige une évaluation et repère la capacité à reprendre.}\end{docu}

\Page{Chapitre 07 / Programme officiel}{Filtrage linéaire}{chapitre-07}
Un filtre transforme l’amplitude et la phase de chaque composante d’un signal. L’outil essentiel est la linéarité : connaître la réponse à une sinusoïde permet de reconstruire la réponse à une somme de sinusoïdes. Le calcul d’une fonction de transfert sert donc à prévoir une sortie et à choisir un dispositif.
\begin{loi}{Notions et capacités travaillées}
Valeurs moyenne et efficace ; décomposition harmonique fournie ; fonctions de transfert d’ordre 1 et 2 ; Bode du premier ordre ; filtres passe-bas, passe-haut et passe-bande ; moyenneur, intégrateur, dérivateur ; cascade et impédances ; exemple mécanique ; simulation spectrale ; détection de non-linéarités.
\end{loi}

\Sub{Prérequis}
Chapitre 6 : amplitudes complexes, impédances, gain, phase et résonance. Logarithmes décimaux ; moyenne d’un cosinus ; dérivation et intégration.

\begin{methode}{La méthode du chapitre}
Repérer les composantes de l’entrée ; calculer ou lire H pour chacune ; multiplier les amplitudes par le gain et ajouter les phases ; reconstruire la sortie ; vérifier sa moyenne, son spectre et sa valeur efficace.
\end{methode}

\source{Concordance : annexe 2 de l’arrêté du 5 janvier 2021, rubrique 1.5. Les séries de Fourier utilisées sont fournies. Le tracé systématique exigible de Bode porte sur l’ordre 1 ; les transferts d’ordre 2 sont exploités et interprétés.}

\Page{Diagnostic / 15 min}{Avant d’ouvrir le cours}{diagnostic}
\begin{questions}
\item Que vaut la moyenne de $\cos\omega t$ sur une période ?
\item $s=2+3\cos\omega t$ V : quelle moyenne ?
\item Convertir un gain 0,1 en décibels.
\item Que signifie une décade ?
\item Une entrée contient 100 Hz et 300 Hz. Un filtre linéaire stable peut-il créer 200 Hz en régime établi ?
\item Que devient C en limite de basse fréquence ?
\item Le module du produit de deux complexes ?
\item Rappeler le transfert RC, sortie C.
\end{questions}
\begin{docu}{Bilan personnel}
\checkline{Les conversions et les unités sont maîtrisées.}\checkline{Je peux orienter un schéma et annoncer les hypothèses.}\checkline{Je sais expliquer une condition initiale.}Corrigé p.\,\pageref{corr-diag}.
\end{docu}


\Page{Cours / 1}{Décrire un signal périodique}{cours-1}
\begin{loi}{Une décomposition fournie}
\[e(t)=a_0+\sum_{n=1}^{N}A_n\cos(n\omega_1t+\varphi_n).\]$a_0$ est la composante continue ; $f_1=\omega_1/(2\pi)$ la fréquence fondamentale ; $nf_1$ la fréquence de l’harmonique n. Les $A_n$ sont des amplitudes positives, les signes étant absorbés dans les phases.
\end{loi}
\Sub{Spectre et forme temporelle}
Un spectre d’amplitude représente des raies à leurs fréquences ; il ne montre pas à lui seul les phases. Deux signaux ayant les mêmes raies peuvent avoir des formes temporelles différentes. Une raie absente a une amplitude nulle ; la fondamentale peut être absente même si la période commune reste définie par les composantes présentes.
\Sub{Exemple}
$e=1+3\cos\omega_1t+\cos(3\omega_1t-\pi/2)$ V a une moyenne de 1 V et deux composantes variables de fréquences $f_1$ et $3f_1$. La période est $1/f_1$. Le spectre à 0 Hz doit être distingué d’une sinusoïde de fréquence très faible.
\begin{methode}{Lire avant de calculer}
Faire un tableau : fréquence, amplitude, phase. Conserver ce tableau pour appliquer H à chaque ligne ; traiter la composante continue par H(0).
\end{methode}


\Page{Cours / 2}{Moyenne et valeur efficace}{cours-2}
\begin{loi}{Définitions}
\[\langle s\rangle=\frac1T\int_0^T s(t)\dd t,\qquad S_{\rm eff}=\sqrt{\frac1T\int_0^T s^2(t)\dd t}.\]Pour $s=S_m\cos(\omega t+\varphi)$, l’identité $\cos^2u=(1+\cos2u)/2$ donne $S_{\rm eff}=S_m/\sqrt2$.
\end{loi}
\Sub{Orthogonalité des harmoniques}
La moyenne du produit de deux cosinus de fréquences harmoniques distinctes est nulle. Les termes croisés disparaissent donc dans la moyenne du carré :\[S_{\rm eff}^2=a_0^2+\sum_{n=1}^N\frac{A_n^2}{2}.\]La contribution de la composante continue doit être conservée. La valeur efficace ne se calcule pas en additionnant les amplitudes, ni les valeurs efficaces.
\Sub{Exemple résolu}
Pour $s=2+3\cos\omega t+\cos3\omega t$ V, $\langle s\rangle=2$ V et $S_{\rm eff}=\sqrt{4+9/2+1/2}=3$ V. Dans une résistance R, la puissance moyenne est $S_{\rm eff}^2/R$ si s est une tension.
\begin{piege}{Quel mode du multimètre ?}
Un mode AC peut supprimer la composante continue. Un appareil donnant une valeur efficace vraie dans sa bande n’a pas le même comportement qu’un appareil étalonné seulement pour des sinusoïdes. Annoncer la grandeur mesurée.
\end{piege}


\Page{Cours / 3}{La fonction de transfert}{cours-3}
\begin{loi}{Un rapport complexe}
\[H(j\omega)=\frac{\underline S}{\underline E},\quad G(\omega)=|H|,\quad \varphi(\omega)=\arg H.\]Pour un filtre linéaire invariant dans le temps, en régime établi :\[s(t)=H(0)a_0+\sum_n A_nG(n\omega_1)\cos[n\omega_1t+\varphi_n+\varphi(n\omega_1)].\]Pour un système réel, H(0) est réel ; son signe est conservé.
\end{loi}
\fig{spectre}{Exemple de passe-bas avec $f_c=f_1$. Les raies de sortie sont légèrement décalées horizontalement pour la lisibilité ; leurs fréquences physiques sont inchangées.}\Sub{La linéarité est le point central}
Le filtre agit séparément sur chaque raie. Une sinusoïde seule reste sinusoïdale, à la même fréquence. La sortie d’un signal composé peut changer de forme parce que ses harmoniques sont transmises différemment.


\Page{Cours / 4}{Passe-bas du premier ordre}{cours-4}
\fig{rc-filtre}{Le condensateur reçoit la tension de sortie. Le modèle suppose une forte impédance de mesure.}\begin{loi}{Réponse harmonique}
\[H=\frac1{1+j\omega/\omega_c},\quad \omega_c=\frac1{RC}.\]\[G=\frac1{\sqrt{1+x^2}},\quad\varphi=-\arctan x,\quad x=\omega/\omega_c.\]À x = 1, $G=1/\sqrt2$ et $\varphi=-\pi/4$.
\end{loi}
\Sub{Comprendre les limites}
À basse fréquence, le condensateur est ouvert et s suit e : H tend vers 1. À haute fréquence, le condensateur dérive le signal et H tend vers zéro. Le gain continu n’est pas nécessairement 1 pour tous les filtres passe-bas ; ici il vaut 1 parce que le montage est non chargé sous source idéale.
\Sub{Exemple}
R = 10 k$\Omega$ et C = 10 nF donnent $\omega_c=10^4$ rad/s et $f_c=1591{,}55$ Hz. À f = fc, une entrée de 2 V d’amplitude donne 1,414 V en sortie, avec retard d’un huitième de période.


\Page{Cours / 5}{Bode : lire une échelle logarithmique}{cours-5}
\begin{loi}{Gain en décibels}
\[G_{\rm dB}=20\log_{10}G.\]Pour le passe-bas d’ordre 1 : $G_{\rm dB}=-10\log_{10}(1+x^2)$. À $x\ll1$, l’asymptote vaut 0 dB ; à $x\gg1$, elle vaut $-20\log_{10}x$, pente -20 dB par décade.
\end{loi}
\fig{bode-bas}{Bode en gain du passe-bas. Le point de coupure est à -3,0103 dB ; -3 dB est la notation usuelle arrondie.}\Sub{La phase fait partie du diagramme}
Elle tend vers 0 à basse fréquence, vaut $-45^\circ$ à la coupure et tend vers $-90^\circ$ à haute fréquence. Sur un tracé, noter les unités de l’axe horizontal : une pulsation ne se lit pas comme une fréquence. Les asymptotes donnent une estimation, pas la valeur exacte près du raccord.
\begin{methode}{Exploiter une pente}
Un décalage de deux décades dans la zone de pente -20 dB par décade donne -40 dB supplémentaires, donc un rapport d’amplitudes 0,01.
\end{methode}


\Page{Cours / 6}{Passe-haut et dérivation}{cours-6}
\fig{cr-filtre}{Le même couple R, C devient passe-haut si la sortie est prise sur R.}\begin{loi}{Transfert et limites}
\[H=\frac{jx}{1+jx},\quad G=\frac{x}{\sqrt{1+x^2}},\quad\varphi=\frac\pi2-\arctan x.\]À basse fréquence H $\simeq j\omega RC$ : $s\simeq RC\dot e$. À haute fréquence H $\simeq1$. La dérivation exige que toutes les composantes utiles vérifient $\omega RC\ll1$.
\end{loi}
\fig{bode-haut}{Le passe-haut gagne +20 dB par décade à basse fréquence, puis atteint le plateau de 0 dB.}\begin{piege}{Une dérivation dans une bande}
Un créneau contient des harmoniques élevées : le modèle de dérivateur ne peut pas rester valable pour toutes. Les pics observés ont une largeur limitée par la bande du montage et de l’appareil.
\end{piege}


\Page{Cours / 7}{Moyenner et intégrer}{cours-7}
\Sub{Le passe-bas comme moyenneur}
Pour conserver la composante continue et atténuer toutes les composantes variables, choisir $f_c\ll f_1$. Alors $s\simeq a_0$. Le rapport entre ondulation résiduelle et moyenne dépend aussi des amplitudes des harmoniques ; la seule condition qualitative sur fc ne fixe pas une précision.
\begin{loi}{Le passe-bas comme intégrateur}
Si chaque composante variable utile vérifie $\omega RC\gg1$, alors\[H\simeq\frac1{j\omega RC},\qquad s_{\rm variable}\simeq\frac1{RC}\int e_{\rm variable}(t)\dd t.\]La composante continue est transmise avec le gain H(0), elle n’est pas intégrée indéfiniment. L’approximation ne décrit pas le démarrage ni une dérive sans borne.
\end{loi}
\Sub{Calculer une contrainte}
Exiger $G(f_1)\leq0{,}01$ pour un passe-bas donne $f_1/f_c\geq\sqrt{9999}$, soit $f_c\leq f_1/99{,}995$. À $f_1=1$ kHz, $f_c\leq10{,}0005$ Hz. Avec C = 1 $\mu\mathrm F$, cela impose $R\geq15{,}915$ k$\Omega$.
\begin{methode}{Nommer l’usage}
Un même passe-bas peut lisser une tension avec moyenne non nulle et transformer sa faible ondulation en intégrale approchée. Préciser la composante concernée.
\end{methode}


\Page{Cours / 8}{Filtres du deuxième ordre}{cours-8}
\begin{loi}{Deux modèles à exploiter}
Passe-bas :\[H_{\rm PB}=\frac{K}{1-x^2+jx/Q}.\]Passe-bande :\[H_{\rm BP}=\frac{K}{1+jQ(x-1/x)}.\]Pour K réel positif, le premier transmet le continu et décroît comme $K/x^2$ à haute fréquence. Le second annule les très basses et très hautes fréquences, avec gain K à x = 1.
\end{loi}
\Sub{Résonance et pentes}
Le passe-bas peut présenter une bosse si $Q>1/\sqrt2$ ; il décroît asymptotiquement de 40 dB par décade. Le passe-bande a des pentes +20 puis -20 dB par décade et une largeur $\Delta\omega=\omega_0/Q$. Q détermine la sélectivité, pas seulement l’amplitude d’une bosse.
\Sub{Choisir selon une tâche}
Pour conserver une moyenne : passe-bas. Pour supprimer un offset lent tout en gardant une bande plus élevée : passe-haut. Pour isoler une composante autour d’une fréquence connue : passe-bande, avec une bande assez large pour garder les variations utiles. Un pic très aigu peut déformer un signal modulé en rejetant ses fréquences voisines.
\begin{docu}{Le calcul sert la lecture}
Le programme demande d’exploiter des transferts d’ordre 2 donnés ou des graphes. Un tracé complet systématique de tous leurs cas ne constitue pas l’objectif de cette partie.
\end{docu}


\Page{Cours / 9}{Cascades et impédances}{cours-9}
\fig{cascade}{Le produit des transferts exige que les conditions de charge restent celles des transferts annoncés.}\Sub{Deux cellules RC non isolées}
Pour deux cellules R, C identiques directement en cascade, la deuxième consomme du courant au nœud de la première. Avec $p=j\omega RC$, la loi des nœuds donne\[H=\frac1{1+3p+p^2}.\]Le produit $(1+p)^{-2}$ correspondrait à deux cellules isolées, par exemple par un suiveur idéal. Les deux expressions ont le même ordre à haute fréquence, mais des gains et phases différents autour de la transition.
\begin{methode}{Condition pratique}
La tension livrée par un étage de sortie de Thévenin $Z_s$ à une entrée $Z_e$ est multipliée par $Z_e/(Z_s+Z_e)$. Pour limiter la perturbation, rechercher $|Z_e|\gg|Z_s|$ sur toute la bande utile. Une impédance est complexe et dépend de la fréquence.
\end{methode}
\Sub{Exemple au raccord}
À $\omega RC=1$, la cascade isolée a H = $-j/2$, gain 0,50 ; la cascade directe a H = $-j/3$, gain 0,333. Multiplier les transferts à vide surestimerait ici la sortie de 50 \%.


\Page{Cours / 10}{Un filtre peut être mécanique}{cours-10}
\Sub{Masse embarquée sur un support mobile}
Un bâti a une position y(t), une masse m une position absolue x(t), et le déplacement relatif mesuré est z = x - y. Le ressort et l’amortisseur exercent $-kz-b\dot z$. La loi de Newton donne\[m\ddot z+b\dot z+kz=-m\ddot y.\]Le transfert de déplacement relatif vaut\[\frac{\underline Z}{\underline Y}=\frac{x_f^2}{1-x_f^2+jx_f/Q},\quad x_f=\omega/\omega_0.\]
\begin{loi}{Deux usages du même dispositif}
À $\omega\ll\omega_0$, $z\simeq-\ddot y/\omega_0^2$ : le déplacement relatif renseigne sur l’accélération du bâti. À $\omega\gg\omega_0$, $\underline Z/\underline Y\simeq-1$ : la masse reste presque immobile tandis que le bâti se déplace, usage sismométrique dans ce modèle.
\end{loi}
\Sub{Interpréter le filtre}
Ce transfert supprime les déplacements très lents du support et transmet les variations rapides avec inversion de signe ; il peut résonner près de sa fréquence propre. La bande utile dépend de la grandeur que l’on souhaite mesurer. Ne pas confondre transfert déplacement-déplacement et transfert accélération-déplacement.
\begin{piege}{Les unités contrôlent le choix}
Z/Y est sans dimension ; Z/A a l’unité s². Un même graphe ne représente pas ces deux transferts.
\end{piege}


\Page{Cours / 11}{Repérer une non-linéarité}{cours-11}
\Sub{Ce que conserve un système linéaire}
Un système linéaire invariant dans le temps ne crée pas de fréquence absente de l’entrée en régime établi. Il peut supprimer des composantes, changer leurs phases et leurs amplitudes ; la déformation temporelle résultante ne prouve donc pas à elle seule une non-linéarité.
\begin{loi}{Un test simple}
Si $s=ae+be^2$ et $e=A\cos\omega t$, alors\[s=aA\cos\omega t+\frac{bA^2}{2}[1+\cos2\omega t].\]L’apparition de continu et de 2f à partir d’un sinus pur révèle ici la non-linéarité. La composante harmonique varie comme $A^2$.
\end{loi}
\Sub{Protocole critique}
Contrôler d’abord le spectre réel de la source. Attendre la disparition du transitoire. Vérifier la fréquence d’échantillonnage et la fenêtre d’analyse ; une fuite spectrale ou un repliement peut simuler des raies. Refaire avec plusieurs amplitudes et vérifier la proportionnalité de la sortie.
\begin{methode}{Conclure avec précision}
Un sinus écrêté donne des harmoniques : une saturation est probable. L’absence de nouvelle raie dans une seule expérience ne prouve pas la linéarité pour toutes les amplitudes et fréquences. La conclusion reste limitée à la plage testée.
\end{methode}


\Page{Pratique / TP}{Filtrer un signal composé}{tp}
\fig{rc-filtre}{Mesurer simultanément e et s, avec les masses au même point.}\Sub{Protocole}
Choisir R = 10 k$\Omega$, C = 100 nF, donc fc = 159,15 Hz. Mesurer d’abord le gain et la phase pour f/fc = 0,1 ; 1 ; 10, avec amplitude d’entrée constante. Puis appliquer un créneau centré et comparer les spectres avant et après filtrage. Garder les mêmes réglages d’analyse pour comparer les raies.
\Sub{Exploitation corrigée}
Les gains attendus aux trois points sont 0,9950 ; 0,7071 ; 0,09950 ; phases $-5{,}71^\circ$, $-45^\circ$, $-84{,}29^\circ$. Pour une composante à 3fc, le gain vaut $1/\sqrt{10}=0{,}3162$. Chaque amplitude de sortie doit être comparée à sa propre raie d’entrée. Le filtre ne change aucune fréquence.
\begin{docu}{Compte rendu}
Présenter le montage orienté, les composants mesurés, les limites instrumentales, les deux diagrammes de Bode et un tableau spectral. Si l’entrée possède un offset, utiliser le mode DC pour étudier son passage. Les valeurs annoncées ici sont des prévisions du modèle, pas des mesures réalisées.
\end{docu}


\Page{Pratique / Numérique}{Reconstruire une sortie à partir du spectre}{numerique}
Le spectre fourni contient une moyenne de 1 V et trois harmoniques. Le code réalise exactement la capacité numérique du programme : simuler l’action du filtre, puis modifier sa coupure.
\begin{verbatim}
import numpy as np
import matplotlib.pyplot as plt
f1 = 100.
t = np.linspace(0, 3/f1, 6001)
n = np.array([1, 3, 5])
A = np.array([3., 1., 0.6])
phi = np.array([0., 0., 0.])
w = 2*np.pi*f1*n
e = 1 + np.sum(A[:,None]*np.cos(
    w[:,None]*t + phi[:,None]), axis=0)
outputs = {}
for fc in (10., 100., 1000.):
    H = 1/(1 + 1j*n*f1/fc)
    s = 1 + np.sum((A*np.abs(H))[:,None]*np.cos(
        w[:,None]*t + (phi+np.angle(H))[:,None]),
        axis=0)
    outputs[fc] = s
    plt.plot(t, s, label=f"fc={fc} Hz")
plt.plot(t, e, color="grey", label="entree")
plt.xlabel("t (s)"); plt.ylabel("tension (V)")
plt.legend(); plt.grid(); plt.tight_layout()
\end{verbatim}
À fc = 100 Hz, les amplitudes de sortie sont 2,1213 ; 0,3162 ; 0,1177 V. La moyenne reste 1 V. À 10 Hz, la sortie s’approche de la moyenne ; à 1000 Hz, elle conserve mieux les détails. Vérifier aussi les phases : un filtrage spectral ne consiste pas seulement à réduire les amplitudes.


\Page{Bilan de cours}{Les réflexes à conserver}{synthese}
\begin{loi}{Les relations et leurs hypothèses}
$S_{\rm eff}^2=a_0^2+\sum A_n^2/2$. Chaque composante est multipliée par H à sa fréquence. $G_{\rm dB}=20\log_{10}|H|$. Pour RC : $\omega_c=1/(RC)$. Les approximations intégrateur et dérivateur exigent une bande adaptée.
\end{loi}
\Sub{Une résolution rédigée}
Partir du besoin et du spectre, puis choisir la nature et la bande du filtre. Vérifier les impédances lorsqu’on raccorde des étages. Reconstruire les phases et contrôler le signal complet.
\begin{piege}{Erreurs à éviter}
Multiplier une série temporelle entière par un gain unique ; additionner les valeurs efficaces ; confondre fréquence et rang harmonique ; négliger la composante continue ; multiplier des transferts à vide malgré une charge significative.
\end{piege}
\begin{docu}{S’autoévaluer}
\checkline{Je retrouve les relations à partir du modèle.}\checkline{Je sais interpréter les limites d’une expression.}\checkline{Je peux dessiner et commenter les courbes.}\checkline{Je relie une mesure à la grandeur calculée.}
\end{docu}


\Page{Exercices / Applications}{Moyenne et valeur efficace}{td1}
\source{Énoncés OptiTAB originaux. Les durées sont indicatives ; justifier toute réponse.}
\ex{E01}{Une sinusoïde décalée}{Niveau 1}{10 à 15 min}
$s=3+4\cos\omega t$ V. Calculer moyenne et valeur efficace.
\ex{E02}{Deux harmoniques}{Niveau 1}{10 à 15 min}
$s=2\cos\omega t+\cos(3\omega t+\pi/4)$ V. Donner la valeur efficace.
\ex{E03}{Un créneau fourni}{Niveau 1}{10 à 15 min}
Un signal vaut +A pendant T/2 puis -A pendant T/2. Utiliser les définitions intégrales pour sa moyenne et sa valeur efficace.


\Page{Exercices / Applications}{Fréquences et décibels}{td2}
\source{Énoncés OptiTAB originaux. Les durées sont indicatives ; justifier toute réponse.}
\ex{E04}{Un spectre incomplet}{Niveau 1}{10 à 15 min}
Un signal contient uniquement des raies à 200 et 300 Hz, de phases fixes et d’amplitudes non nulles. Donner sa période minimale.
\ex{E05}{Gains et décibels}{Niveau 1}{10 à 15 min}
Convertir G = 2, G = 0,01 et $G_{\rm dB}=-6$ dB.
\ex{E06}{Une pente}{Niveau 1}{10 à 15 min}
Un gain vaut -10 dB à 1 kHz puis suit une pente de -40 dB par décade. Estimer le gain à 100 kHz.


\Page{Exercices / Applications}{Étudier un passe-bas}{td3}
\source{Énoncés OptiTAB originaux. Les durées sont indicatives ; justifier toute réponse.}
\ex{E07}{Choisir C}{Niveau 2}{10 à 15 min}
R = 10 k$\Omega$ et fc = 1,00 kHz pour un RC. Déterminer C.
\ex{E08}{Une composante à la coupure}{Niveau 2}{10 à 15 min}
Un passe-bas RC reçoit $e=4\cos(2\pi f_ct)$ V. Écrire s.
\ex{E09}{Au-dessus de la coupure}{Niveau 2}{10 à 15 min}
Pour un passe-bas, calculer gain et phase à f = 10fc. Comparer à l’asymptote intégratrice.


\Page{Exercices / Méthodes}{Étudier un passe-haut}{td4}
\source{Énoncés OptiTAB originaux. Les durées sont indicatives ; justifier toute réponse.}
\ex{E10}{Une moyenne supprimée}{Niveau 2}{10 à 15 min}
Un passe-haut reçoit $e=2+3\cos\omega_ct$ V. Déterminer s en régime établi.
\ex{E11}{Dériver une composante}{Niveau 2}{10 à 15 min}
Pour un passe-haut avec $\omega RC=0{,}02$, donner le signal approché si $e=5\cos\omega t$ V.
\ex{E12}{Retrouver la coupure}{Niveau 2}{10 à 15 min}
Un passe-haut présente une phase de $45^\circ$ à 800 Hz. Quel est fc ? Que vaut son gain ?


\Page{Exercices / Méthodes}{Reconstruire la sortie}{td5}
\source{Énoncés OptiTAB originaux. Les durées sont indicatives ; justifier toute réponse.}
\ex{E13}{Filtrer trois raies}{Niveau 2}{15 à 25 min}
$e=1+2\cos\omega_ct+\cos3\omega_ct$ V entre dans un passe-bas RC. Écrire s et sa valeur efficace.
\ex{E14}{La phase manquante}{Niveau 2}{15 à 25 min}
Deux filtres ont le même module H à toutes les harmoniques utiles, mais des phases différentes. Leurs sorties ont-elles même valeur efficace et même forme ?
\ex{E15}{Éliminer un bruit}{Niveau 2}{15 à 25 min}
Une composante à 50 kHz doit être atténuée d’au moins 40 dB par un passe-bas d’ordre 1. Donner fc maximal.


\Page{Exercices / Méthodes}{Choisir un filtre}{td6}
\source{Énoncés OptiTAB originaux. Les durées sont indicatives ; justifier toute réponse.}
\ex{E16}{Un moyenneur précis}{Niveau 2}{15 à 25 min}
$e=2+\cos(2\pi1000t)$ V. Exiger une amplitude d’ondulation inférieure à 0,02 V. Trouver fc maximal d’un passe-bas.
\ex{E17}{Une bande sélective}{Niveau 2}{15 à 25 min}
Un passe-bande canonique possède f0 = 2 kHz, Q = 20, K = 3. Donner gain maximal et largeur.
\ex{E18}{Une seconde pente}{Niveau 2}{15 à 25 min}
Pour $H=1/(1-x^2+jx/Q)$ avec Q = 0,50, trouver G à x = 1 et la limite à grand x.


\Page{Exercices / Approfondissement}{Cascades et mécanique}{td7}
\source{Énoncés OptiTAB originaux. Les durées sont indicatives ; justifier toute réponse.}
\ex{E19}{Deux cellules isolées}{Niveau 3}{15 à 25 min}
Deux passe-bas RC identiques sont isolés par un suiveur idéal. À $\omega RC=1$, déterminer le transfert total.
\ex{E20}{Deux cellules raccordées}{Niveau 3}{15 à 25 min}
Pour deux RC identiques sans suiveur, utiliser $H=1/(1+3p+p^2)$ avec p = j\ensuremath{\omega}RC. Calculer H à \ensuremath{\omega}RC = 1 et comparer à E19.
\ex{E21}{Une masse sur un bâti}{Niveau 3}{15 à 25 min}
Le déplacement relatif vérifie $\ddot z+2\dot z+100z=-\ddot y$. Donner \ensuremath{\omega_0} et Q ; interpréter le cas \ensuremath{\omega} très faible.


\Page{Exercices / Approfondissement}{Tester la linéarité}{td8}
\source{Énoncés OptiTAB originaux. Les durées sont indicatives ; justifier toute réponse.}
\ex{E22}{Un étage trop chargé}{Niveau 3}{15 à 25 min}
Un étage a une sortie résistive 1 k$\Omega$ et l’entrée suivante vaut 9 k$\Omega$. Quel facteur de tension supplémentaire apparaît ?
\ex{E23}{Un carré révélateur}{Niveau 3}{15 à 25 min}
$s=e+0{,}2e^2/\mathrm V$ avec $e=2\cos\omega t$ V. Donner le spectre de sortie.
\ex{E24}{Un test insuffisant}{Niveau 3}{15 à 25 min}
Un signal composé est déformé en sortie mais aucune fréquence nouvelle n’apparaît. Peut-on conclure à une non-linéarité ? Proposer un meilleur test.


\Page{Défi avancé / D1}{Conserver un signal, rejeter une perturbation}{defi1}
\label{ex-D1}\source{Problème OptiTAB original. Niveau 3 ; 40 à 60 minutes. Indices p.\,\pageref{indices} ; corrigé p.\,\pageref{corr-defi1}.}
Un passe-bas d’ordre 1 doit transmettre une composante utile à 100 Hz avec gain au moins 0,95 et atténuer une perturbation à 10 kHz avec gain au plus 0,02.\begin{questions}
\item Traduire les deux conditions en bornes sur fc. Le choix est-il possible ?
\item Choisir fc = 400 Hz et R = 10 k$\Omega$ ; déterminer C.
\item Pour $e=1+2\cos(2\pi100t)+0{,}5\cos(2\pi10000t)$ V, écrire la sortie et calculer sa valeur efficace.
\item Quelle composante fixe l’ondulation résiduelle ? Pourquoi une sortie proche du signal utile ne constitue-t-elle pas un moyenneur ?
\end{questions}


\Page{Défi avancé / D2}{Retrouver le transfert d’une cascade réelle}{defi2}
\label{ex-D2}\source{Problème OptiTAB original. Niveau 3 ; 40 à 60 minutes. Indices p.\,\pageref{indices} ; corrigé p.\,\pageref{corr-defi2}.}
Deux résistances R sont en série entre l’entrée e et la sortie s. Un condensateur C relie le nœud intermédiaire u à la masse ; un autre C relie la sortie à la masse. La sortie est mesurée à vide.\begin{questions}
\item Écrire les deux lois des nœuds en amplitudes complexes.
\item Éliminer u et établir H en posant p = j\ensuremath{\omega}RC.
\item Identifier \ensuremath{\omega_0} et Q dans la forme canonique du deuxième ordre.
\item Comparer le continu, le gain à \ensuremath{\omega}RC = 1 et les hautes fréquences avec deux RC isolés.
\item Pourquoi observe-t-on la même pente finale mais pas la même fréquence de transition ?
\end{questions}


\Page{Défi avancé / D3}{Un accéléromètre mécanique}{defi3}
\label{ex-D3}\source{Problème OptiTAB original. Niveau 3 ; 40 à 60 minutes. Indices p.\,\pageref{indices} ; corrigé p.\,\pageref{corr-defi3}.}
Le déplacement relatif z d’une masse sur un support vérifie $\ddot z+(\omega_0/Q)\dot z+\omega_0^2z=-a(t)$, où a est l’accélération du support. On choisit Q = $1/\sqrt2$ et \ensuremath{\omega_0} = 1000 rad/s.\begin{questions}
\item Établir le transfert Z/A et sa limite statique.
\item On estime l’accélération par $a_{\rm est}=-\omega_0^2z$. Calculer le rapport complexe entre estimation et accélération.
\item Calculer l’erreur sur le module et la phase à \ensuremath{\omega} = 100 rad/s.
\item Pourquoi un bon accord en amplitude n’assure-t-il pas une bonne reproduction temporelle ?
\item Que devient ce capteur à fréquence très élevée ?
\end{questions}


\Page{Évaluation / DM}{DM — Extraire une moyenne}{dm}
\source{Devoir maison OptiTAB original ; 1 h 30 à 2 h ; calculatrice autorisée. Barème pédagogique sur 20.}
Une tension vaut $e=2+3\cos(2\pi1000t)+\cos(2\pi3000t)$ V. On utilise un passe-bas RC avec R = 16 k$\Omega$ et C = 1,0 $\mu\mathrm F$.\Sub{A. Spectre et efficacité — 5 points}
\begin{questions}
\item Donner moyenne, période et valeur efficace de l’entrée. (3)
\item Pourquoi le spectre d’amplitude ne suffit-il pas à tracer tout signal ? (2)
\end{questions}

\Sub{B. Sortie — 9 points}
\begin{questions}
\item Établir H et calculer fc. (3)
\item Donner amplitudes et phases des deux composantes variables de sortie. (4)
\item Calculer moyenne et valeur efficace de sortie. (2)
\end{questions}

\Sub{C. Cahier des charges — 6 points}
\begin{questions}
\item L’amplitude d’ondulation doit rester sous 0,04 V. Donner une borne supérieure simple à partir des deux amplitudes et conclure. (3)
\item Expliquer ce qu’afficheraient un voltmètre DC, un mode AC idéal et un oscilloscope en couplage DC. (3)
\end{questions}



\Page{Évaluation / DS / 1 sur 2}{DS — Séparer des composantes}{ds1}
\source{Devoir surveillé OptiTAB original ; durée 2 h ; calculatrice autorisée ; note sur 20. Les deux pages forment un seul sujet.}
\Sub{A. Un signal à traiter — 5 points}
$e=1+2\cos(2\pi200t)+\cos(2\pi2000t)$ V.\begin{questions}
\item Donner son spectre, sa période et sa valeur efficace. (3)
\item Quel type de filtre permet de réduire la composante à 2 kHz en gardant 200 Hz ? Quelles données manquent pour choisir sa coupure ? (2)
\end{questions}

\Sub{B. Un premier ordre — 9 points}
On choisit $H=1/(1+jf/f_c)$, fc = 400 Hz.\begin{questions}
\item Justifier les deux asymptotes de Bode et donner gain et phase à fc. (3)
\item Reconstruire s(t), puis calculer sa moyenne et sa valeur efficace. (4)
\item Ce filtre est-il un intégrateur pour la composante à 200 Hz ? (2)
\end{questions}



\Page{Évaluation / DS / 2 sur 2}{Approfondir et conclure}{ds2}
\Sub{C. Un second ordre donné — 6 points}
On remplace le filtre par $H_2=1/[1-(f/f_0)^2+j(f/f_0)/Q]$, avec f0 = 400 Hz et $Q=1/\sqrt2$.\begin{questions}
\item Montrer que $|H_2|=1/\sqrt{1+(f/f_0)^4}$. (2)
\item Comparer les gains à 200 Hz et 2 kHz avec le premier ordre. (2)
\item Une raie à 4 kHz apparaît en sortie mais n’existe pas à l’entrée. Donner une conclusion prudente et deux vérifications expérimentales. (2)
\end{questions}

\begin{docu}{Gestion du temps}
Lecture : 10 min ; parties A et B : 70 min ; partie C : 30 min ; vérification : 10 min. Une expression littérale, une unité et une interprétation accompagnent les résultats.
\end{docu}


\Page{Aides progressives}{Des indices pour repartir}{indices}
\Sub{Défi D1}
1. Isoler fc dans le gain exact. 2. Garder les phases. 3. La valeur efficace contient le continu.
\Sub{Défi D2}
1. Au nœud de sortie : u = (1+p)s. 2. Au premier nœud : e = (2+p)u-s. 3. Développer avant de remplacer p par j\ensuremath{\omega}RC.
\Sub{Défi D3}
1. Le signe moins doit être conservé. 2. L’estimation suit un passe-bas de gain statique 1. 3. Avec ce Q, le module se simplifie.
\Sub{Exercices}
E04 : une fondamentale peut être absente. E13 : filtrer séparément les raies. E15 et E16 : utiliser le gain exact avant une approximation. E20 : le transfert direct n’est pas le produit à vide.
\Sub{Évaluations}
Pour borner une ondulation, la somme des amplitudes donne une borne sûre. Une valeur efficace utilise les carrés. Pour un cahier des charges, établir chaque inégalité avant de choisir un composant.


\Page{Corrigés / Diagnostic}{Réactiver les bons réflexes}{corr-diag}
\Sub{D1}
Zéro ; l’intégrale de son carré vaut T/2.

\Sub{D2}
2 V. L’amplitude 3 V ne modifie pas la moyenne.

\Sub{D3}
$20\log_{10}(0{,}1)=-20$ dB pour un rapport d’amplitudes.

\Sub{D4}
Un rapport 10 de fréquences : par exemple de 100 Hz à 1 kHz.

\Sub{D5}
Non, pour un système linéaire invariant dans le temps. Une composante initiale transitoire doit être éteinte.

\Sub{D6}
Un circuit ouvert dans le modèle idéal.

\Sub{D7}
Le produit des modules ; l’argument est la somme des arguments modulo $2\pi$.

\Sub{D8}
$H=1/(1+jRC\omega)$ à vide, sous source idéale.


\Page{Corrigés / Exercices}{Moyenne et valeur efficace}{corr-td1}
\sol{E01}{Une sinusoïde décalée}
$\langle s\rangle=3$ V. $S_{\rm eff}=\sqrt{9+16/2}=\sqrt{17}=4{,}123$ V. Le terme croisé a une moyenne nulle.
\sol{E02}{Deux harmoniques}
Les composantes de fréquences distinctes sont orthogonales : $S_{\rm eff}=\sqrt{2^2/2+1^2/2}=\sqrt{2{,}5}=1{,}581$ V, indépendante ici de la phase relative.
\sol{E03}{Un créneau fourni}
Les deux contributions à l’intégrale de s se compensent : moyenne nulle. $s^2=A^2$ presque partout, donc valeur efficace A. La relation $A/\sqrt2$ est réservée à une sinusoïde.


\Page{Corrigés / Exercices}{Fréquences et décibels}{corr-td2}
\sol{E04}{Un spectre incomplet}
Les fréquences sont 2 et 3 fois 100 Hz. La période commune minimale est 10 ms, même si la raie à la fondamentale de 100 Hz est absente. Il faut chercher le plus grand diviseur commun des fréquences.
\sol{E05}{Gains et décibels}
$20\log_{10}2=6{,}0206$ dB ; 0,01 donne -40 dB. À -6 dB, G = $10^{-6/20}=0{,}5012$.
\sol{E06}{Une pente}
Deux décades donnent -80 dB supplémentaires : total -90 dB. Le gain d’amplitude est $10^{-4{,}5}=3{,}162\times10^{-5}$. La réponse est asymptotique, si la pente annoncée reste valable.


\Page{Corrigés / Exercices}{Étudier un passe-bas}{corr-td3}
\sol{E07}{Choisir C}
$C=1/(2\pi Rf_c)=15{,}915$ nF. Ne pas utiliser $1/(Rf_c)$ : la coupure en pulsation vaut $1/(RC)$.
\sol{E08}{Une composante à la coupure}
$s=2\sqrt2\cos(2\pi f_ct-\pi/4)$ V. Le gain est $1/\sqrt2$ et le retard un huitième de période.
\sol{E09}{Au-dessus de la coupure}
$G=1/\sqrt{101}=0{,}09950$ et phase $-84{,}289^\circ$. L’approximation donne G = 0,10 et phase $-90^\circ$. L’erreur relative d’amplitude est environ 0,50 \%, mais l’écart de phase reste 5,71°.


\Page{Corrigés / Exercices}{Étudier un passe-haut}{corr-td4}
\sol{E10}{Une moyenne supprimée}
H(0) = 0 ; la moyenne disparaît. À la coupure, $H=(1+j)/2$, gain $1/\sqrt2$, phase $+\pi/4$. Donc $s=(3/\sqrt2)\cos(\omega_ct+\pi/4)$ V.
\sol{E11}{Dériver une composante}
$s\simeq RC\dot e=-0{,}10\sin\omega t=0{,}10\cos(\omega t+\pi/2)$ V. Le module exact est $5(0{,}02)/\sqrt{1+0{,}02^2}=0{,}09998$ V.
\sol{E12}{Retrouver la coupure}
Pour l’ordre 1 canonique, $\varphi=\pi/2-\arctan(f/f_c)$. La phase $\pi/4$ impose f/fc = 1. Donc fc = 800 Hz et G = $1/\sqrt2$.


\Page{Corrigés / Exercices}{Reconstruire la sortie}{corr-td5}
\sol{E13}{Filtrer trois raies}
$s=1+\sqrt2\cos(\omega_ct-\pi/4)+(1/\sqrt{10})\cos(3\omega_ct-\arctan3)$ V. $S_{\rm eff}^2=1+2/2+0{,}1/2=2{,}05$ ; $S_{\rm eff}=1{,}432$ V.
\sol{E14}{La phase manquante}
Même valeur efficace pour le même spectre d’entrée, puisque seuls les carrés des amplitudes interviennent. La forme temporelle peut différer par les phases relatives. Si toutes les phases supplémentaires correspondent à un même retard, la forme est simplement translatée.
\sol{E15}{Éliminer un bruit}
$G\leq0{,}01$ impose $(50000/f_c)^2\geq9999$, donc $f_c\leq50000/\sqrt{9999}=500{,}025$ Hz. Vérifier ensuite que les composantes utiles restent transmises.


\Page{Corrigés / Exercices}{Choisir un filtre}{corr-td6}
\sol{E16}{Un moyenneur précis}
$1/\sqrt{1+(1000/f_c)^2}\leq0{,}02$. Donc $f_c\leq1000/\sqrt{2499}=20{,}004$ Hz. La moyenne de sortie vaut 2 V et la limite porte sur l’amplitude de l’ondulation.
\sol{E17}{Une bande sélective}
Le maximum vaut K = 3, soit 9,542 dB ; $\Delta f=f_0/Q=100$ Hz. Un signal utile s’étendant sur 500 Hz serait fortement déformé par ce choix.
\sol{E18}{Une seconde pente}
À x = 1, G = Q = 0,50. À grand x, $H\simeq-1/x^2$, donc G $\simeq1/x^2$, phase tendant vers $-\pi$ et pente -40 dB par décade. Il n’y a pas de bosse de résonance.


\Page{Corrigés / Exercices}{Cascades et mécanique}{corr-td7}
\sol{E19}{Deux cellules isolées}
$H=(1+j)^{-2}=1/(2j)=-j/2$. Gain 0,50, phase $-\pi/2$. Les gains en décibels s’additionnent : deux fois -3,0103 dB.
\sol{E20}{Deux cellules raccordées}
$H=1/[1+3j-1]=1/(3j)=-j/3$. Gain 1/3 au lieu de 1/2 ; la phase reste $-\pi/2$ à ce point. La deuxième cellule charge la première.
\sol{E21}{Une masse sur un bâti}
$\omega_0=10$ rad/s et $\omega_0/Q=2$, donc Q = 5. Pour \ensuremath{\omega} très faible, $z\simeq-\ddot y/100$ : la position relative mesure l’accélération du bâti avec le facteur $-0{,}01$ s².


\Page{Corrigés / Exercices}{Tester la linéarité}{corr-td8}
\sol{E22}{Un étage trop chargé}
$9/(1+9)=0{,}90$, soit -0,915 dB. Une perte de 10 \% peut être significative ; le produit des transferts à vide doit être corrigé par ce facteur si le modèle résistif est valable.
\sol{E23}{Un carré révélateur}
$s=2\cos\omega t+0{,}4+0{,}4\cos2\omega t$ V. Les nouvelles composantes sont 0 et 2f ; le système est non linéaire. Le coefficient 0,2 par volt assure l’homogénéité.
\sol{E24}{Un test insuffisant}
Non : une modification différente des amplitudes et phases suffit à déformer un signal dans un filtre linéaire. Injecter un sinus pur, vérifier le spectre d’entrée, rechercher des harmoniques et faire varier l’amplitude pour tester la proportionnalité, en évitant saturation et artefacts d’acquisition.


\Page{Corrigé du défi D1}{Conserver un signal, rejeter une perturbation}{corr-defi1}
À 100 Hz, $f_c\geq100/\sqrt{1/0{,}95^2-1}=304{,}24$ Hz. À 10 kHz, $f_c\leq10000/\sqrt{2499}=200{,}04$ Hz. Les deux bornes sont incompatibles : un premier ordre ne satisfait pas ce cahier des charges. Le choix fc = 400 Hz sert à quantifier l’écart, il ne le résout pas. C = 39,789 nF. La sortie vaut $1+(2/\sqrt{1+0{,}25^2})\cos(2\pi100t-\arctan0{,}25)+(0{,}5/\sqrt{626})\cos(2\pi10000t-\arctan25)$ V. Amplitudes : 1,9403 V et 0,019984 V ; phases -14,036° et -87,709°. $S_{\rm eff}=\sqrt{1+1{,}9403^2/2+0{,}019984^2/2}=1{,}6978$ V. Le gain du bruit vaut 0,03997, supérieur à 0,02 malgré sa faible amplitude initiale. Le signal à 100 Hz est conservé ; la sortie n’est donc pas une moyenne. Un ordre plus élevé ou une autre architecture doit être étudié avec le même cahier des charges.

\Page{Corrigé du défi D2}{Retrouver le transfert d’une cascade réelle}{corr-defi2}
Les lois sont $(\underline E-\underline U)/R=\underline UjC\omega+(\underline U-\underline S)/R$ et $(\underline U-\underline S)/R=jC\omega\underline S$. La seconde donne U = (1+p)S ; la première donne E = (2+p)U-S. Ainsi E = $(1+3p+p^2)$S et $H=1/(1+3p+p^2)$. Avec x = \ensuremath{\omega}RC, le dénominateur est $1-x^2+3jx$ : $\omega_0=1/(RC)$ et Q = 1/3. À continu, les deux montages ont gain 1. À x = 1, le montage direct a gain 1/3 ; l’isolé a gain 1/2 et Q = 1/2 dans sa forme canonique. À grand x, les deux tendent vers $-1/x^2$. L’ordre fixe la pente finale, mais le coefficient du terme du premier degré règle l’amortissement et la transition. Une pente seule ne permet pas d’identifier tout un transfert.

\Page{Corrigé du défi D3}{Un accéléromètre mécanique}{corr-defi3}
$\underline Z/\underline A=-1/[\omega_0^2-\omega^2+j\omega\omega_0/Q]$. La limite statique est $-1/\omega_0^2=-10^{-6}$ s². Le rapport $\underline A_{\rm est}/\underline A=1/(1-x^2+j\sqrt2x)$. Son module vaut $1/\sqrt{1+x^4}$. À x = 0,1, le gain est 0,999950, soit une erreur relative d’amplitude d’environ $5\times10^{-5}$ ; la phase vaut $-\arctan(\sqrt2\times0{,}1/0{,}99)=-8{,}13^\circ$. La quasi-perfection du module masque un retard mesurable. Pour plusieurs harmoniques, leurs phases différentes peuvent déformer le signal. À grand x, le rapport estimé tend vers zéro comme $1/x^2$ et sa phase vers $-\pi$ ; le capteur ne mesure plus fidèlement l’accélération. Une calibration doit préciser une bande de module et une tolérance de phase.

\Page{Corrigé / DM / Sur 20}{DM — Extraire une moyenne}{corr-dm}
\Sub{A — 5 points}
Moyenne 2 V ; période 1 ms ; $E_{\rm eff}=\sqrt{4+9/2+1/2}=3$ V (3). Les phases relatives sont nécessaires à la reconstruction temporelle (2).
\Sub{B — 9 points}
$H=1/(1+jRC\omega)$ et fc = 9,9472 Hz (3). À 1000 Hz, x = 100,531, amplitude $3/\sqrt{1+x^2}=0{,}029840$ V, phase -89,430°. À 3000 Hz, amplitude 0,0033157 V, phase -89,810° (4). La moyenne reste 2 V ; $S_{\rm eff}=\sqrt{4+(0{,}029840^2+0{,}0033157^2)/2}=2{,}000113$ V (2).
\Sub{C — 6 points}
L’amplitude absolue de l’ondulation est bornée par la somme des amplitudes, soit 0,033156 V < 0,04 V ; la contrainte est satisfaite sans rechercher le maximum exact (3). Le DC mesure la moyenne 2 V ; le mode AC efficace idéal mesure l’ondulation efficace $\sqrt{(0{,}029840^2+0{,}0033157^2)/2}=0{,}02123$ V. L’oscilloscope DC montre la faible ondulation autour de 2 V (3).


\Page{Corrigé / DS / 1 sur 2}{Calculer et interpréter}{corr-ds1}
\Sub{A — 5 points}
Raies : continu 1 V, 200 Hz amplitude 2 V, 2000 Hz amplitude 1 V. Période 5 ms. $E_{\rm eff}=\sqrt{1+2+0{,}5}=\sqrt{3{,}5}$ V (3). Passe-bas ; il faut préciser les tolérances de transmission à 200 Hz et de rejet à 2 kHz (2).
\Sub{B — 9 points}
À bas f : H tend vers 1, 0 dB. À haut f : H $\simeq f_c/(jf)$, pente -20 dB/décade. À fc : $G=1/\sqrt2$, phase -45° (3). La sortie vaut $1+(2/\sqrt{1{,}25})\cos(2\pi200t-\arctan0{,}5)+(1/\sqrt{26})\cos(2\pi2000t-\arctan5)$ V. Amplitudes 1,7889 et 0,19612 V, phases -26,565° et -78,690°. Moyenne 1 V ; $S_{\rm eff}=\sqrt{1+1{,}6+1/52}=1{,}6184$ V (4). À 200 Hz, f/fc = 0,5 ne vérifie pas la condition de haute fréquence : approximation intégratrice injustifiée (2).


\Page{Corrigé / DS / 2 sur 2}{Calculer et interpréter}{corr-ds2}
\Sub{C — 6 points}
Le dénominateur au carré vaut $(1-x^2)^2+2x^2=1+x^4$ (2). À x = 0,5, $G_2=4/\sqrt{17}=0{,}97014$, supérieur à 0,89443. À x = 5, $G_2=1/\sqrt{626}=0{,}039968$, inférieur à 0,19612 : meilleur maintien de 200 Hz et meilleur rejet de 2 kHz (2). Une fréquence nouvelle est incompatible avec le modèle linéaire invariant en régime établi. Vérifier le spectre source, la saturation, l’échantillonnage et les transitoires avant de conclure à une non-linéarité du filtre (2).


\Page{Support authentique / 1}{TD publié — Deux filtres réels}{sources-1}
\source{Source consultée le 25 octobre 2026. TD7, 2025–2026 : exercice 1 et exercice 8.}
\begin{loi}{Accéder au document d’origine}
\href{https://mpsi.schleck.ovh/fichiers/td/td7_filtrage_lineaire.pdf}{Ouvrir le sujet original et conserver ses figures et données.}
\end{loi}
\Sub{Parcours de travail}
Traiter le filtre RL sortie bobine, puis la cascade RC directe. Redessiner les ports et les masses.
\Sub{Repères de correction indépendants}
Pour RL : $H=jL\omega/(R+jL\omega)$, passe-haut de coupure R/L. Pour deux RC identiques directement raccordés : $H=1/[1+3jRC\omega-(RC\omega)^2]$. Le second étage modifie le premier ; la simple multiplication à vide échoue.
\source{Ce parcours ne reproduit pas l’épreuve complète. Les questions non ciblées peuvent demander des chapitres ultérieurs. Le temps de travail proposé concerne uniquement cette sélection.}

\Page{Support authentique / 2}{DS publié — Le signal d’une guitare}{sources-2}
\source{Source consultée le 25 octobre 2026. DS n°4 du 11 décembre 2025, partie I, p. 1–2. Sélection : I.1.1 puis I.1.2 a et b.}
\begin{loi}{Accéder au document d’origine}
\href{https://mpsi.schleck.ovh/fichiers/devoirs/2025-2026/ds04_enonce.pdf}{Ouvrir le sujet original et conserver ses figures et données.}
\end{loi}
\Sub{Parcours de travail}
Lire le plateau de basse fréquence et la pente de haute fréquence. Choisir parmi les quatre transferts proposés et établir les asymptotes, sans reprendre les parties de chimie.
\Sub{Repères de correction indépendants}
Le modèle est un passe-bas d’ordre 2 : fonction 2, $H=H_0/(1-x^2+jx/Q)$. À basse fréquence, $G_{\rm dB}\simeq20\log_{10}H_0$ ; à haute fréquence, $G_{\rm dB}\simeq20\log_{10}H_0-40\log_{10}x$, avec $x=\omega/\omega_0$. À $\omega_0$, $H=-jH_0Q$ : gain $H_0Q$, phase $-\pi/2$. Les valeurs de H0 et Q se déterminent ensuite avec les niveaux du graphe, question c.
\source{Ce parcours ne reproduit pas l’épreuve complète. Les questions non ciblées peuvent demander des chapitres ultérieurs. Le temps de travail proposé concerne uniquement cette sélection.}

\Page{Références et suivi}{Préparer la suite}{references}
\Sub{Référence réglementaire}
\href{https://www.education.gouv.fr/bo/21/Special1/ESRS2035779A.htm}{Bulletin officiel spécial n° 1 du 11 février 2021, arrêté du 5 janvier 2021.}\par
\href{https://cache.media.education.gouv.fr/file/SPE1-MEN-MESRI-4-2-2021/64/8/spe779_annexe_1373648.pdf}{Annexe officielle regroupée}, rubrique 1.5. Les séries de Fourier utilisées sont fournies. Le tracé systématique exigible de Bode porte sur l’ordre 1 ; les transferts d’ordre 2 sont exploités et interprétés.
\Sub{Nature des contenus}
Le cours, les 24 exercices, les trois défis et les évaluations OptiTAB sont rédigés pour cette collection. Les graphiques sont calculés avec les équations annoncées. Les supports authentiques sont identifiés sur les pages précédentes ; leurs énoncés et figures restent dans le document d’origine.
\Sub{Compétences à confirmer}
\checkline{Je peux retrouver l’équation sans mémoriser une suite de signes.}
\checkline{Je vérifie les conditions initiales et les limites.}
\checkline{Je peux confronter un calcul, un graphe et un bilan énergétique.}
\checkline{Je sais corriger mon DS et refaire une question sans aide.}
\begin{methode}{Prochaine étape}
Le chapitre 8 étudie un autre effet sur un signal : le retard dû à sa propagation dans l’espace.
\end{methode}

